\diarytitle{0104}{$\arctan(2x/(1-x^2))$ の導関数}

合成関数の微分公式 $\dfrac{d}{dx}\arctan u = \dfrac{u'}{1+u^{2}}$（連鎖律）を用いる。

$u = \dfrac{2x}{1-x^{2}}$ とおく。商の微分法より
\[
u' = \frac{2(1-x^{2}) - 2x\cdot(-2x)}{(1-x^{2})^{2}}
= \frac{2 - 2x^{2} + 4x^{2}}{(1-x^{2})^{2}}
= \frac{2 + 2x^{2}}{(1-x^{2})^{2}}
= \frac{2(1+x^{2})}{(1-x^{2})^{2}}
\]
また
\[
1 + u^{2} = 1 + \frac{4x^{2}}{(1-x^{2})^{2}}
= \frac{(1-x^{2})^{2} + 4x^{2}}{(1-x^{2})^{2}}
= \frac{1 - 2x^{2} + x^{4} + 4x^{2}}{(1-x^{2})^{2}}
= \frac{(1+x^{2})^{2}}{(1-x^{2})^{2}}
\]
したがって
\[
\frac{dy}{dx} = \frac{\dfrac{2(1+x^{2})}{(1-x^{2})^{2}}}{\dfrac{(1+x^{2})^{2}}{(1-x^{2})^{2}}}
= \frac{2(1+x^{2})}{(1+x^{2})^{2}}
= \frac{2}{1+x^{2}}
\]
（$x^{2} \neq 1$、すなわち $x \neq \pm 1$ で定義される。）したがって
\[
\boxed{\; \frac{dy}{dx} = \frac{2}{1+x^{2}} \;}
\]
（検算: これは $\arctan(2x/(1-x^{2})) = 2\arctan x$（$|x|<1$、正接の2倍角公式 $\tan(2\theta) = \dfrac{2\tan\theta}{1-\tan^{2}\theta}$ に $\theta=\arctan x$ を代入した恒等式）と一致することからも確認できる。実際 $y=2\arctan x$ を直接微分すると $\dfrac{dy}{dx} = 2\cdot\dfrac{1}{1+x^{2}} = \dfrac{2}{1+x^{2}}$ となり、上の結果と一致する。）
